% L3a data plate: daily continuously compounded growth rates of SPY (top) and
% AMD (bottom), 2014--2024, computed from daily closes with Delta t = 1/252 yr.
% These are the series the stylized-facts example notebook works with (tickers
% from MyTrainingMarketDataSet in VLQuantitativeFinancePackage).
%
% Each day is a vertical bar from zero (ycomb), not a connected line: with
% 2766 observations a line reads as a fuzzy band, while bars render the
% volatility ENVELOPE, which is what clustering is about. The x axis is
% labeled by calendar year (tick positions are the first trading day of each
% year in the t index).
%
% The stacked layout carries the lesson. SPY makes clustering unmistakable:
% its rolling 21-day std is ~15/yr inside the marked March-2020 band against
% ~0.6/yr in the calmest 2017 window. The same band is turbulent in AMD, so
% volatility clusters are market-wide weeks, not one ticker's quirk. The
% April-2016 AMD spike is the counterexample: ONE large jump, not a cluster.
%
% SIZED FOR THE DECK: the slide includes this PDF at natural size, so keep the
% total height near 5cm and every label INSIDE its panel (a node above the
% axis would grow the bounding box).
%
% Data files: Data-L3a-{SPY,AMD}-DailyGrowthRate.csv (t index, date, g in
% 1/yr), regenerated from the course dataset; `date` is unused by the plot.
\documentclass[tikz,border=2pt]{standalone}
\usepackage{vnfigure}
\usepackage{pgfplots}
\usepgfplotslibrary{groupplots}
\pgfplotsset{compat=1.18}

% Type roles, matching the equity-market plate's convention.
\newcommand{\fAxis}{\sffamily\fontsize{8.2}{10}\selectfont}   % ticks, axis labels
\newcommand{\fNote}{\sffamily\fontsize{9}{11}\selectfont}     % in-panel annotations

\begin{document}
\begin{tikzpicture}
  \begin{groupplot}[
      group style={
        group size=1 by 2,
        vertical sep=0.3cm,
        xticklabels at=edge bottom,
      },
      width=14.2cm, height=3.35cm,
      axis lines=left,
      axis line style={vnink, line width=0.6pt},
      tick style={vnink, line width=0.4pt},
      ticklabel style={font=\fAxis, color=vnink},
      ylabel style={color=vnink},
      xmin=-20, xmax=2800,
      xtick={1,503,1006,1509,2014,2515},
      xticklabels={2014,2016,2018,2020,2022,2024},
  ]
    % ---- SPY: the index shows clean calm/turbulent regimes ----
    \nextgroupplot[
        ymin=-33, ymax=26,
        ytick={-20,0,20},
        ylabel={\fAxis SPY $g$ (1/yr)},
    ]
      \fill[vncarnelian!10] (axis cs:1540,-33) rectangle (axis cs:1600,26);
      \draw[vncarnelian!60, dashed, line width=0.6pt] (axis cs:1540,-33) -- (axis cs:1540,26);
      \draw[vncarnelian!60, dashed, line width=0.6pt] (axis cs:1600,-33) -- (axis cs:1600,26);
      \node[anchor=west, text=vnmuted, font=\fNote] at (axis cs:1650,19)
        {a high-volatility cluster};
      \addplot [ycomb, vnlink, line width=0.3pt] table [x=t, y=g, col sep=comma]
        {Data-L3a-SPY-DailyGrowthRate.csv};

    % ---- AMD: same turbulent weeks, plus a single-jump counterexample ----
    \nextgroupplot[
        ymin=-80, ymax=112,
        ytick={-50,0,50,100},
        ylabel={\fAxis AMD $g$ (1/yr)},
    ]
      \fill[vncarnelian!10] (axis cs:1540,-80) rectangle (axis cs:1600,112);
      \draw[vncarnelian!60, dashed, line width=0.6pt] (axis cs:1540,-80) -- (axis cs:1540,112);
      \draw[vncarnelian!60, dashed, line width=0.6pt] (axis cs:1600,-80) -- (axis cs:1600,112);
      \node[anchor=west, text=vnmuted, font=\fNote] at (axis cs:660,88)
        {a single large jump};
      \addplot [ycomb, vnlink, line width=0.3pt] table [x=t, y=g, col sep=comma]
        {Data-L3a-AMD-DailyGrowthRate.csv};
  \end{groupplot}
\end{tikzpicture}
\end{document}
